\section{UVOD}
Intralogistika obuhvaća organizaciju i upravljanje unutarnjim tokom materijala i robe. Zbog rasta e-trgovine, zahtjeva za bržom isporukom i nedostatka radne snage, tradicionalni sustavi postaju usko grlo. Klasični viličari često su neekonomični i podložni ljudskim pogreškama.

Povijest transportnih sustava bez vozača počinje 1953. kada je Arthur Mac Barrett predstavio prvo \cite{barrett} automatski vođeno vozilo (AGV) koje je pratilo žicu u podu. Značajan iskorak napravio je Volvo 1973. godine uvođenjem flote AGV-ova u proizvodnju automobila, čime je povećana fleksibilnost. Kasnije su se razvile navigacijske metode poput magnetskih traka i laserskog vođenja (LGV), no sustavi su i dalje ovisili o infrastrukturi.

Početkom 21. stoljeća razvoj senzora i umjetne inteligencije doveo je do autonomnih mobilnih robota (AMR) i autonomnih viličara (AGF). Oni koriste algoritme poput SLAM-a za lokalizaciju i navigaciju bez potrebe za vanjskom infrastrukturom. U okviru Industrije 4.0, ta vozila postaju inteligentni agenti koji se integriraju sa sustavima za upravljanje skladištem (WMS) i planiranjem (ERP). S razvojem koncepta Industrije 5.0, fokus se pomiče na suradnju čovjeka i stroja (\textit{human-centric} pristup). U tom kontekstu razvijaju se autonomni viličari s mogućnošću dvostrukog načina rada (hibridni viličari). Takva vozila obavljaju rutinske zadatke potpuno autonomno, no opremljena su i standardnim komandama koje operateru omogućuju ručno upravljanje za obavljanje složenijih, nestandardnih ili hitnih zadataka, čime čovjek ostaje ključni, neizostavni dio sustava.
Cilj ovog rada je istražiti razvoj i primjenu autonomnih viličara. Rad analizira razlike između klasičnih i autonomnih viličara, sigurnosne sustave (ISO 3691-4), navigacijske tehnologije i praktične primjere iz industrije, te daje pregled prednosti i ograničenja uvođenja ove tehnologije.

% TODO: Dodan kvantitativni motiv
Prema podacima \cite{dzs,ecommerce_report}, skladišni incidenti povezani s radom viličara čine 15\% svih ozljeda u logistici. Rast elektroničke trgovine od 20\% godišnje povećava zahtjeve za operativnom dostupnošću i točnošću isporuke; zato automatizacija intralogistike postaje ključna kako bi se osigurale 24/7 operacije i smanjile pogreške u inventaru.
